\medskip
3. \underline{Cohomology of quadrics}

3.1. Let $p:X\longrightarrow S$ be a smooth quadric of dimension $n$
over $S$, and let $\ell$ be a prime number invertible on $S$.  We write
\begin{equation*}
\eta\in H^0\bigl(S,R^2p_*\mathbb Z_\ell(1)\bigr)
\end{equation*}
for the ``cohomology class of a hyperplane section.''  Locally on $S$
for the \'{e}tale topology, $P(X)\simeq\mathbb P(V^*)$; by definition,
$\eta$ is then the image in
$H^0(S,R^2p_*\mathbb Z_\ell(1))$ of the first Chern class
\begin{equation*}
c_1(\mathcal O(1))\in H^2(X,\mathbb Z_\ell(1)).
\end{equation*}
\footnote{The displayed definition in the print has $R^1p_*$, whereas
the following sentence and the degree of $c_1(\mathcal O(1))$ both give
$R^2p_*$.}

3.2. Suppose that $n=2m$.  A generator $D$ of $X$ then defines a
cohomology class
$c\ell'(D)\in H^n(X,\mathbb Z_\ell(m))$, whose image in
$H^0(S,R^np_*\mathbb Z_\ell(m))$ is denoted $c\ell(D)$.  By 2.8, the
latter depends only on the section $e(D)$ of $Z(X)$.  We therefore obtain
a morphism of \'{e}tale sheaves on $S$
\begin{equation}\tag{3.2.1}
c\ell:Z(X)\longrightarrow R^np_*\mathbb Z_\ell(m).
\end{equation}

\underline{Theorem} 3.3. \underline{Let $p:X\longrightarrow S$ be a
smooth quadric of dimension $n$ over $S$, and let $\ell$ be a prime
number invertible on $S$.}

\begin{enumerate}[label=(\roman*),leftmargin=2.6em]
\item $R^{2i+1}p_*\mathbb Z_\ell=0$.

%%%%%%%%%% source scan idx 83 | folio 76 %%%%%%%%%%

\item If $0\leq 2i<n$ (respectively, if $n<2i\leq 2n$), then
$R^{2i}p_*\mathbb Z_\ell(i)$ is canonically isomorphic to the constant
sheaf $\mathbb Z_\ell$, and is generated by $\eta^i$ (respectively, by
$\eta^i/2$).

\item If $n=2m$, the map deduced from (3.2.1),
\begin{equation}\tag{3.3.1}
\mathbb Z_\ell^{Z(X)}\longrightarrow R^np_*\mathbb Z_\ell(m),
\end{equation}
is an isomorphism.  If $\alpha$ and $\beta$ are two disjoint sections
of $Z(X)$ (assumed isomorphic to $S\amalg S$, as is locally the case),
then
\begin{align*}
\text{(a)}\quad &\eta^m=c\ell(\alpha)+c\ell(\beta),\\[1ex]
\text{(b)}\quad &
\begin{cases}
\operatorname{Tr}(c\ell(\alpha)^2)
=\operatorname{Tr}(c\ell(\beta)^2)=1,
& c\ell(\alpha)c\ell(\beta)=0,\quad m\text{ even},\\[1ex]
c\ell(\alpha)^2=c\ell(\beta)^2=0,
& \operatorname{Tr}(c\ell(\alpha)c\ell(\beta))=1,
  \quad m\text{ odd},
\end{cases}\\[1ex]
\text{(c)}\quad &\eta\bigl(c\ell(\alpha)-c\ell(\beta)\bigr)=0.
\end{align*}
\end{enumerate}

\footnote{The print repeats the tag (3.2.1) on the map in part (iii),
but the proof itself refers to this map as (3.3.1).}

Assertions (i) and (ii) are recalled only for reference (XI, 1.6,
supplemented by XI, 2.6 when $2i+1=n$ in (i)).

It is enough to verify (iii) when $S$ is the spectrum of an
algebraically closed field $k$ (one could even take
$S=\Spec(\mathbb C)$).

\underline{Formula (a).}  We may suppose that $X$ is the quadric in
$\mathbb P^{2m+1}(k)$ with equation
\begin{equation*}
\sum_{i=0}^{m}x_i x_{i+m+1}=0.
\end{equation*}
The class $\eta^m$ is the cohomology class of the plane section defined
by
\begin{equation*}
x_i=0\qquad(0\leq i<m),
\end{equation*}
and this cycle is the sum of the generators $D_1,D_2$ defined by

%%%%%%%%%% source scan idx 84 | folio 77 %%%%%%%%%%

\begin{align*}
D_1:&\quad x_i=0 &&(0\leq i\leq m),\\
D_2:&\quad x_i=0 &&(0\leq i<m),\qquad x_{2m+1}=0.
\end{align*}
By 1.12, $e(D_1)\neq e(D_2)$, and hence
\begin{equation*}
\eta^m=c\ell(D_1)+c\ell(D_2)
=c\ell(\alpha)+c\ell(\beta).
\end{equation*}

\footnote{The printed ambient projective space is
$\mathbb P^{m+1}$, although the displayed nondegenerate equation has
$2m+2$ coordinates.  The required ambient space is
$\mathbb P^{2m+1}$.}

\underline{Formula (b).}  If generators $D_1,D_2$ are disjoint
(respectively, meet transversally in one point), then
$c\ell(D_1)c\ell(D_2)=0$ (respectively,
$\operatorname{Tr}(c\ell(D_1)c\ell(D_2))=1$).  By 1.12, when $m$ is
even we have $e(D_1)\neq e(D_2)$ (respectively,
$e(D_1)=e(D_2)$); when $m$ is odd we have $e(D_1)=e(D_2)$
(respectively, $e(D_1)\neq e(D_2)$).  Formula (b) follows.

\underline{Formula (c).}  This is clear when $n=0$.  Otherwise, by
(ii), it is enough to prove
\begin{equation*}
\eta^m\bigl(c\ell(\alpha)-c\ell(\beta)\bigr)=0.
\end{equation*}
Indeed,
\begin{equation*}
\operatorname{Tr}\!\left(
\eta^m(c\ell(\alpha)-c\ell(\beta))\right)
=\operatorname{Tr}\!\left(c\ell(\alpha)^2-c\ell(\beta)^2\right)=0.
\end{equation*}

\underline{Verification.}  When $m$ is even,
\begin{equation*}
2=\operatorname{Tr}(\eta^{2m})
=\operatorname{Tr}\!\left(c\ell(\alpha)^2
+2c\ell(\alpha)c\ell(\beta)+c\ell(\beta)^2\right)=1+1;
\end{equation*}
when $m$ is odd,
\begin{equation*}
2=\operatorname{Tr}(\eta^{2m})
=\operatorname{Tr}\!\left(c\ell(\alpha)^2
+2c\ell(\alpha)c\ell(\beta)+c\ell(\beta)^2\right)=2.
\end{equation*}
\footnote{In these two top-degree intersection computations the print
has $\eta^2$ in place of $\eta^{2m}$, and once prints
$c\ell(\alpha^2)$ in place of $c\ell(\alpha)^2$.}

That (3.3.1) is an isomorphism now follows from the fact that the two
bilinear forms in two variables
\begin{equation*}
B(x,y)=x_1y_2+x_2y_1,
\qquad
B(x,y)=x_1y_1+x_2y_2
\end{equation*}
have discriminant $\pm1$.

%%%%%%%%%% source scan idx 85 | folio 78 %%%%%%%%%%

\underline{Verification} 3.4. Let $X$ be a smooth quadric of dimension
$n$ over a finite field $\mathbb F_q$.  In view of 3.3, the Lefschetz
trace formula gives the classical formula
\begin{equation*}
\#X(\mathbb F_q)=
\begin{cases}
\displaystyle\sum_{i=0}^{n}q^i,
  & n\text{ odd},\\[2ex]
\displaystyle\sum_{i=0}^{n}q^i+\varepsilon q^m,
  & n=2m,\quad \varepsilon=\pm1.
\end{cases}
\end{equation*}
When $n$ is even, $\varepsilon=1$ if $X$ is split, that is, if it has a
rational generator, and $\varepsilon=-1$ otherwise.

3.5. Let $f:X\longrightarrow S$ be a smooth quadric of dimension
$n=2m$.  The \underline{primitive part} of
$R^nf_*\mathbb Z_\ell(m)$ is the orthogonal complement of $\eta^m$;
with the notation of 3.3, it is generated by
$c\ell(\alpha)-c\ell(\beta)$.  The \underline{primitive quotient} of
$R^nf_*\mathbb Z_\ell(m)$ is the quotient by
$\mathbb Z_\ell\eta^m$.  The primitive part and primitive quotient of
$R^nf_*\mathbb Z_\ell$ are obtained by tensoring with
$\mathbb Z_\ell(-m)$.

